% -------------------
%   Abstrakt - Slovensky
% -------------------
\newpage 
\section*{Abstrakt}

Centrálne banky, ako je Európska centrálna banka, neurčujú iba úrokové sadzby --- používajú aj jazyk na to, aby usmerňovali očakávania finančných trhov o budúcnosti. Táto práca skúma, či slová samotnej Európskej centrálnej banky dokážu predpovedať vývoj úrokových sadzieb a reakcie trhov. Pomocou metód spracovania prirodzeného jazyka sme prečítali všetky tlačové konferencie Európskej centrálnej banky z rokov 1998 až 2024 a emocionálny tón každej pasáže sme zmerali dvoma špecializovanými jazykovými modelmi. Metodologickou inováciou je rozdelenie každej tlačovej konferencie na dve časti: vopred pripravený prejav, ktorý predstavitelia banky starostlivo nacvičia, a spontánny rozhovor, v ktorom novinári kladú nepripravené otázky. Zistili sme, že novinári sú systematicky pesimistickejší ako samotní centrálni bankári a že ignorovanie tohto rozdielu skresľuje signál o menovej politike v starších štúdiách, ktoré obe časti spájajú do jedného dokumentu. Keď oba hlasy oddelíme, emocionálny tón centrálnej banky funguje ako systém včasného varovania: modely strojového učenia natrénované na tomto textovom signáli dokážu niekoľko mesiacov vopred predpovedať, či nasledujúce rozhodnutie o úrokovej sadzbe bude zvýšenie, zníženie alebo ponechanie bez zmeny. Na vzorke $289$ tlačových konferencií stúpla krížovo validovaná metrika AUC pri predpovedaní smerovania ďalšieho rozhodnutia ECB o úrokových sadzbách z $0,682$ (iba na základe histórie sadzieb) na $0,763$ (model náhodného lesa využívajúci príznaky sentimentu z modelov FinBERT a CentralBankRoBERTa).

\paragraph*{Kľúčové slová:} 
analýza sentimentu, spracovanie prirodzeného jazyka, menová politika, prediktívna analýza, komunikácia centrálnej banky
% --- Koniec Abstrakt - Slovensky


% -------------------
% --- Abstrakt - Anglicky 
% -------------------
\newpage 
\section*{Abstract}

Modern central banks do more than set interest rates --- they also use language to shape what financial markets expect about the future. This thesis asks whether the words spoken by the European Central Bank can predict where interest rates and market expectations will move. Using natural language processing, we read every European Central Bank press conference from 1998 to 2024 and measured the emotional tone of each passage with two specialised language models. The methodological innovation is to split each press conference into two parts: the prepared speech, which bank officials write and rehearse in advance, and the spontaneous interview that follows it, in which journalists ask unscripted questions. We find that journalists are persistently more pessimistic than the central bankers themselves, and that ignoring this difference contaminates the policy signal in earlier studies that treat the two as a single document. Once the two voices are separated, the emotional tone of the central bank acts as an early-warning system: machine-learning models trained on this textual signal can predict, several months in advance, whether the next interest-rate decision will be a hike, a cut, or no change. Across $289$ press conferences, the cross-validated AUC for predicting the direction of the next ECB rate decision rises from $0.682$ (rate history alone) to $0.763$ (random forest with FinBERT and CentralBankRoBERTa sentiment features).

\paragraph*{Keywords:} 
sentiment analysis, natural language processing, monetary policy, predictive analysis, central bank communication
% --- Koniec Abstrakt - Anglicky